\documentclass[11pt,a4paper]{article}
\usepackage[T1]{fontenc}
\usepackage[utf8]{inputenc}
\usepackage{lmodern}
\usepackage[czech]{babel}
\usepackage[margin=2cm]{geometry}
\usepackage{amsmath,amssymb,mathtools}
\usepackage{xcolor,cancel,enumitem}
\usepackage{tikz,pgfplots}
\pgfplotsset{compat=1.18}
\usetikzlibrary{arrows.meta,calc,patterns,matrix,decorations.pathreplacing,babel}
\pagestyle{empty}
\setlength{\parindent}{0pt}
\setlength{\parskip}{4pt}
\newcommand{\vysl}[1]{\boxed{\color{red!75!black}#1}}
\newcommand{\pozn}[1]{{\color{red!75!black}#1}}
\newcommand{\dx}{\,\mathrm{d}x}
\newcommand{\dt}{\,\mathrm{d}t}
\newcommand{\subst}[1]{\left|\begin{array}{l}#1\end{array}\right|}
\DeclareMathOperator{\arctg}{arctg}
\DeclareMathOperator{\arccotg}{arccotg}
\DeclareMathOperator{\tg}{tg}
\DeclareMathOperator{\cotg}{cotg}
\usepackage{tcolorbox}
\newtcolorbox{pomocny}{colframe=red!75!black, colback=white, coltitle=white, fonttitle=\bfseries, title={Pomocný výpočet:}, sharp corners}
\newcommand{\dy}{\,\mathrm{d}y}
% znaménkové schéma: \znam{bod/popis/typ,...}{znaménka v intervalech}
% typ: o = vyloučený bod (kroužek), x = obyčejný bod
\newcommand{\znam}[2]{%
\begin{tikzpicture}[baseline=-3pt]
\def\n{0}
\foreach \p/\l/\t [count=\i] in {#1} {\xdef\n{\i}}
\pgfmathsetmacro{\L}{1.6*(\n+1)}
\draw (0,0) -- (\L,0);
\foreach \p/\l/\t [count=\i] in {#1} {
  \pgfmathsetmacro{\xx}{1.6*\i}
  \ifnum\pdfstrcmp{\t}{o}=0 \draw[fill=white] (\xx,0) circle (2.5pt); \draw (\xx-0.09,0) -- (\xx+0.09,0) (\xx,-0.09) -- (\xx,0.09);
  \else \draw (\xx,0.1) -- (\xx,-0.1); \fi
  \node[below=2pt] at (\xx,0) {\small $\l$};
}
\foreach \s [count=\i] in {#2} {
  \pgfmathsetmacro{\xx}{1.6*\i-0.8}
  \node[above] at (\xx,0) {$\s$};
}
\end{tikzpicture}}
\newcommand{\aug}[1]{\left(\begin{array}{rrr|rrr}#1\end{array}\right)}
\newcommand{\krok}[1]{\text{\small\color{orange!80!black}#1}}

% Hornerovo schéma jako TikZ matice: řádek 1 = koeficienty, další řádky = kořeny,
% poslední řádek = podílový polynom zapsaný pod koeficienty
\tikzset{hm/.style={matrix of math nodes, nodes in empty cells, column sep=0pt, row sep=0pt,
  nodes={minimum width=2.7em, minimum height=1.55em, anchor=center, inner sep=1pt}}}
% \hornerlines{matice}{poslední řádek schématu}{počet sloupců}{poslední sloupec pod závorkou}
\newcommand{\hornerlines}[4]{%
  \pgfmathtruncatemacro{\hlR}{#2-1}\pgfmathtruncatemacro{\hlN}{#3-1}%
  \foreach \i in {1,...,\hlR} \draw (#1-\i-1.south west) -- (#1-\i-#3.south east);
  \foreach \j in {1,...,\hlN} \draw (#1-1-\j.north east) -- (#1-#2-\j.south east);
  \draw[decorate, decoration={brace, mirror, amplitude=4pt}]
    ([yshift=-2pt]#1-#2-2.south west) -- ([yshift=-2pt]#1-#2-#4.south east);}
\newcommand{\tm}[2]{\tikz[remember picture, baseline=(#1.base)]\node[inner sep=0pt, outer sep=1pt] (#1) {$\displaystyle #2$};}

\newenvironment{hb}{\par\smallskip\centering}{\par\vspace{-1.6\baselineskip}}

\begin{document}

\section*{Kořeny, rozklad, znaménko polynomu}

\textbf{a)}\quad $3x^3-5x^2-4x+4$ \qquad (možné kořeny: $\pm1,\pm2,\pm4$)
\begin{hb}
\begin{tikzpicture}
\matrix (m) [hm] {
  & 3 & -5 & -4 & 4\\
2 & 3 & 1 & -2 & 0\\[0.3cm]
  & 3x^2 & +x & -2 & \\};
\hornerlines{m}{2}{5}{4}
\node[anchor=west] at (m-2-5.east) {$\to\ (x-2)$};
\node[anchor=base west] at (m-3-4.base east) {$=3(x+1)\left(x-\frac23\right)$};
\end{tikzpicture}
\end{hb}
\begin{align*}
D &= 1^2-4\cdot3\cdot(-2)=1+24=25\\
x_{1,2} &= \frac{-1\pm5}{3\cdot2}=\begin{cases}-\frac66=-1\\[2pt] \frac46=\frac23\end{cases}
\end{align*}
rozklad: $3(x-2)(x+1)\left(x-\frac23\right)$\\
kořeny: $x_1=2$, $x_2=-1$, $x_3=\frac23$\\
znaménko: \znam{a/-1/x,b/\frac23/x,c/2/x}{-,+,-,+}

\medskip\hrule\medskip
\textbf{b)}\quad $3x^3-8x^2+7x-2$ \qquad (možné kořeny: $\pm1,\pm2$)
\begin{hb}
\begin{tikzpicture}
\matrix (m) [hm] {
  & 3 & -8 & 7 & -2\\
1 & 3 & -5 & 2 & 0\\[0.3cm]
  & 3x^2 & -5x & +2 & \\};
\hornerlines{m}{2}{5}{4}
\node[anchor=west] at (m-2-5.east) {$\to\ (x-1)$};
\node[anchor=base west] at (m-3-4.base east) {$=3(x-1)\left(x-\frac23\right)$};
\end{tikzpicture}
\end{hb}
\begin{align*}
D &= (-5)^2-4\cdot3\cdot2=25-24=1\\
x_{1,2} &= \frac{5\pm1}{2\cdot3}=\begin{cases}\frac66=1\\[2pt] \frac46=\frac23\end{cases}
\end{align*}
kořeny: $x_1=1$, $x_2=1$, $x_3=\frac23$\\
rozklad: $3(x-1)^2\left(x-\frac23\right)$\\
znaménko:
\begin{tikzpicture}[baseline=-3pt]
\draw (0,0) -- (4.8,0);
\draw (1.6,0.1) -- (1.6,-0.1); \draw (3.2,0.1) -- (3.2,-0.1);
\node[below=2pt] at (1.6,0) {\small $\frac23$};
\node[below=2pt, draw, circle, inner sep=1pt] (r1) at (3.2,0) {\small $1$};
\node[above] at (0.8,0) {$-$}; \node[above] at (2.4,0) {$+$}; \node[above] at (4.0,0) {$+$};
\node[anchor=west] (t1) at ($(r1.east)+(0.8,0)$) {dvojnásobný};
\node[anchor=north west, inner sep=0pt] at ($(t1.south west)+(0.1,0)$) {\small(nemění znaménko)};
\draw[-{Stealth}] (r1.east) -- (t1.west);
\end{tikzpicture}

\newpage
\textbf{c)}\quad $x^4+3x^3-7x^2-9x+12$ \qquad (možné kořeny: $\pm1,\pm2,\pm3,\pm4,\pm6,\pm12$)
\begin{hb}
\begin{tikzpicture}
\matrix (m) [hm] {
   & 1 & 3 & -7 & -9 & 12\\
 1 & 1 & 4 & -3 & -12 & 0\\
-4 & 1 & 0 & -3 & 0 & \\[0.3cm]
   & x^2 & {\color{gray}+0x} & -3 & & \\};
\hornerlines{m}{3}{6}{4}
\node[anchor=west] at (m-2-6.east) {$\to\ (x-1)$};
\node[anchor=west] at (m-3-6.east) {$\to\ (x+4)$};
\node[anchor=base west] at (m-4-4.base east) {$=0$};
\end{tikzpicture}
\end{hb}
\begin{align*}
x^2-3&=0\\
x^2&=3\\
x&=\pm\sqrt3
\end{align*}
kořeny: $x_1=1$, $x_2=-4$, $x_3=\sqrt3$, $x_4=-\sqrt3$\\
rozklad: $(x-1)(x+4)\left(x-\sqrt3\right)\left(x+\sqrt3\right)$\\
znaménko: \znam{a/-4/x,b/-\sqrt3/x,c/1/x,d/\sqrt3/x}{+,-,+,-,+}

\medskip\hrule\medskip
\textbf{d)}\quad $x^4-2x^3+2x-1$ \qquad (možné kořeny: $\pm1$)
\begin{hb}
\begin{tikzpicture}
\matrix (m) [hm] {
  & 1 & -2 & 0 & 2 & -1\\
1 & 1 & -1 & -1 & 1 & 0\\
1 & 1 & 0 & -1 & 0 & \\[0.3cm]
  & x^2 & {\color{gray}+0x} & -1 & & \\};
\hornerlines{m}{3}{6}{4}
\node[anchor=west] at (m-2-6.east) {$\to\ (x-1)$};
\node[anchor=west] at (m-3-6.east) {$\to\ (x-1)$};
\node[anchor=base west] at (m-4-4.base east) {$=0$};
\end{tikzpicture}
\end{hb}
\[ (x-1)(x+1)=0 \]
kořeny: $x_{1,2,3}=1$ (trojnásobný), $x_4=-1$\\
rozklad: $(x-1)^3(x+1)$\\
znaménko: \znam{a/-1/x,b/1/x}{+,-,+}

\newpage
\textbf{e)}\quad $2x^3-3x^2-3x+2$ \qquad (možné kořeny: $\pm1,\pm2$)
\begin{hb}
\begin{tikzpicture}
\matrix (m) [hm] {
   & 2 & -3 & -3 & 2\\
-1 & 2 & -5 & 2 & 0\\[0.3cm]
   & 2x^2 & -5x & +2 & \\};
\hornerlines{m}{2}{5}{4}
\node[anchor=west] at (m-2-5.east) {$\to\ (x+1)$};
\node[anchor=base west] at ($(m-3-4.base east)+(1.2,0)$) {$D=(-5)^2-4\cdot2\cdot2=9$};
\end{tikzpicture}
\end{hb}
\[
x_{1,2}=\frac{5\pm3}{2\cdot2}=\begin{cases}\frac84=2\\[2pt] \frac24=\frac12\end{cases}
\]
kořeny: $x_1=-1$, $x_2=2$, $x_3=\frac12$\\
rozklad: $2(x+1)(x-2)\left(x-\frac12\right)$\\
znaménko: \znam{a/-1/x,b/\frac12/x,c/2/x}{-,+,-,+}

\medskip\hrule\medskip
\textbf{f)}\quad $2x^4-x^3+2x^2-3x = \tm{fx}{x}\left(2x^3-x^2+2x-3\right)$ \qquad (možné kořeny: $\pm1,\pm3$)
\begin{hb}
\begin{tikzpicture}[remember picture]
\matrix (m) [hm] {
  & 2 & -1 & 2 & -3\\
1 & 2 & 1 & 3 & 0\\[0.3cm]
  & 2x^2 & +x & +3 & \\};
\hornerlines{m}{2}{5}{4}
\node[anchor=east] (x2) at ($(m-2-1.west)+(-1.3,0)$) {$x_2=1$};
\node[anchor=east] (x1) at ($(m-1-1.west)+(-1.3,0.05)$) {$x_1=0$};
\draw[-{Stealth}] (m-2-1.west) -- (x2.east);
\end{tikzpicture}
\begin{tikzpicture}[remember picture, overlay]
\draw[-{Stealth}] ([yshift=-1pt]fx.south) to[out=-100, in=80] ($(x1.north west)+(0.2,-0.03)$);
\end{tikzpicture}
\end{hb}
\[ D=1^2-4\cdot2\cdot3<0 \ \text{-- nemá reálné kořeny} \]
kořeny: $x_1=0$, $x_2=1$\\
rozklad: $x(x-1)\left(2x^2+x+3\right)$\\
znaménko: \znam{a/0/x,b/1/x}{+,-,+}

\end{document}
